%% kapa: exam=1 round=30 year=2019 period=2 form=A subject=accounting subject_ko=회계학 seq=2 number=42
%% source: Q-Net 감정평가사 기출문제 (cst00309) artlSeq=5205205 file=extracted/1008/5205205/2019년 제30회 감정평가사 1차시험 문제지/2019년 제30회 감정평가사 1차시험 문제지/2019년도 제30회 감정평가사 1차 2교시 A형.hwp
%% conversion: kordoc 4.24.51 -> tools/gen1.py (auto); review=proofread-sonnet; flags=none
\begin{question}{42}
(주)감평과 (주)한국은 사용 중인 유형자산을 상호 교환하여 취득하였다. 두 회사가 보유하고 있는 유형자산에 대한 자료는 다음과 같으며, 교환 시 (주)감평이 (주)한국에 추가로 현금 ￦200,000을 지급하였다. 이들 자산간 교환취득을 상업적 실질이 있다고 가정할 경우, (주)감평이 인식할 유형자산 취득원가(A)와 (주)한국이 인식할 유형자산처분이익(B)은? (단, 두 자산의 공정가치는 신뢰성 있게 측정할 수 있으며, 각 회사의 입장에서 취득한 자산의 공정가치가 더 명백하다는 증거는 없다.) \par
{\small\centering
\begin{tabular}{|c|c|c|}
\hline
구분 & (주)감평 & (주)한국 \\
\hline
취득원가 & ￦2,250,000 & ￦1,500,000 \\
감가상각누계액 & 1,250,000 & 600,000 \\
공정가치 & 950,000 & 1,150,000 \\
\hline
\end{tabular}\par}

\begin{choices}
  \item A: ￦950,000\hspace{2em}B: ￦250,000
  \item A: ￦950,000\hspace{2em}B: ￦450,000
  \item A: ￦1,050,000\hspace{2em}B: ￦450,000
  \item A: ￦1,150,000\hspace{2em}B: ￦250,000
  \item A: ￦1,150,000\hspace{2em}B: ￦450,000
\end{choices}
\end{question}
\answer{4}
